\section{Discusión preliminar}
La comparación de integradores permitió separar dos conceptos que no deben confundirse: precisión y estabilidad. Con pasos moderados, RK4 alcanzó un error muy inferior al de ambos métodos de Euler, como se aprecia en la tabla \ref{tab:errores} y la figura \ref{fig:convergencia}. El descenso experimental de error de RK4 fue consistente con orden cuatro, mientras los métodos de Euler se comportaron aproximadamente como esquemas de primer orden. No obstante, el rendimiento observado depende del tamaño de paso, de la trayectoria y de la referencia numérica seleccionada. Asimismo, el número de evaluaciones de $\mathbf f$ no representa todo el costo computacional de Euler hacia atrás, que incluye construcción de jacobianas, Newton y Gauss--Seidel.

RKF45 fue especialmente competitivo en evaluaciones para los escenarios de precisión inicial, pero su paso fue adaptativo y sus condiciones no equivalieron a un paso fijo idéntico al de RK4. La corrección de la comparación en los nodos aceptados eliminó una discrepancia artificial de conservación que aparecía cuando se interpolaban externamente las trayectorias. Esa revisión puso de relieve la necesidad de validar no solo el integrador, sino también la forma de medir sus errores. Un estudio posterior debería variar las tolerancias de RKF45 y comparar costo frente a una misma meta de error.

El incremento de la razón espectral local con $1/\varepsilon$ fue compatible con la estructura rápido--lento de Bulai et al. (2024). Los fallos de admisibilidad de los métodos explícitos para pasos grandes, recogidos en la tabla \ref{tab:rigidez} y la figura \ref{fig:rigidez}, mostraron que un paso razonable para la dinámica epidemiológica puede resultar excesivo para la informativa. Sin embargo, el éxito de Euler hacia atrás no se interpretó como estabilidad incondicional del algoritmo completo: Gauss--Seidel puede no converger para ciertas matrices internas y se recurrió a rechazar pasos, verificando positividad y residuos. Estos mecanismos aumentaron las iteraciones y, en algunos casos, el número de pasos efectivo superó ampliamente el esperado con el paso nominal.

Finalmente, el contraste de los escenarios epidemiológicos de la tabla \ref{tab:epidemia} y las figuras \ref{fig:desinformacion}--\ref{fig:escepticismo} mostró variaciones marcadas en la infección asociadas con la presencia de una rama informativa particular. El primer escenario mantuvo una prevalencia apreciable en el horizonte simulado, mientras el control decayó; el segundo produjo el resultado opuesto. Los resultados son compatibles con el acoplamiento propuesto por Bulai et al. (2024), pero dependen de las condiciones iniciales y los parámetros elegidos. El estudio no constituye una estimación epidemiológica real ni una demostración de extinción asintótica. Quedan como limitaciones la ausencia de datos observados para calibración, el carácter local del indicador espectral, el rango finito de pasos explorados y la necesidad de contrastar formalmente la convergencia de las referencias antes de cerrar el paper final.

El barrido uniparamétrico complementó la comparación entre métodos: los valores finales de infección registrados dependieron de $b_1$, $b_2$, $K$ y $\beta$. La Tabla~\ref{tab:parametricos} presenta los extremos explorados; estas variaciones aisladas no permiten extrapolar conclusiones a todos los regímenes ni probar rebrotes no observados.
